%==============================================================
%  lecons/analyse/bernstein.tex — Polynômes de Bernstein
%==============================================================

\lecontitle{Les polynômes de Bernstein}{Analyse réelle — Approximation uniforme}

%=== NOTATION LOCALE =========================================
\newcommand{\Bn}{B_n}
\newcommand{\Bin}{B_{i,n}}

%=============================================================
%  § 1 — DÉFINITION ET RÉSULTATS FONDAMENTAUX
%=============================================================
\secnum{navyblue}{1}{Définition et résultats fondamentaux}

\begin{tcolorbox}[thmbox]
  \textbf{Définition.}\enspace\index{polynômes de Bernstein}
  Pour $n \geq 1$ et $0 \leq i \leq n$, le $i$-ième \emph{polynôme de Bernstein
  de base} de degré $n$ est
  \[
    \Bin(x) \;=\; \binom{n}{i}\,x^i\,{(1-x)}^{n-i}, \qquad x \in [0,1].
  \]
  À toute fonction $f : [0,1] \to \mathbb{R}$, on associe son \emph{$n$-ième
  polynôme de Bernstein}\index{opérateur de Bernstein}
  \[
    \Bn(f)(x) \;=\; \sum_{i=0}^{n} f\!\left(\frac{i}{n}\right)
                    \binom{n}{i}\,x^i\,{(1-x)}^{n-i}.
  \]

  \medskip
  \textbf{Propriétés élémentaires.}
  \begin{itemize}
    \item \textit{Partition de l'unité :}\quad
          $\displaystyle\sum_{i=0}^{n} \Bin(x) = 1$ pour tout $x \in [0,1]$.
    \item \textit{Positivité :}\quad
          $\Bin(x) \geq 0$ sur $[0,1]$, donc $\Bn$ est un opérateur \emph{positif}
          ($f \geq 0 \Rightarrow \Bn(f) \geq 0$).
    \item \textit{Reproduction des fonctions affines :}\quad
          $\Bn(\mathbf{1}) = 1$ et $\Bn(\mathrm{id}) = \mathrm{id}$, autrement dit
          \[
            \Bn(1)(x) = 1, \qquad \Bn(x \mapsto x)(x) = x.
          \]
          Plus généralement, $\Bn$ fixe toute fonction affine.
    \item \textit{Linéarité et monotonie :}\quad $\Bn$ est linéaire en $f$ et
          $f \leq g \Rightarrow \Bn(f) \leq \Bn(g)$.
  \end{itemize}

  \medskip
  \textbf{Théorème (Weierstrass, version constructive).}\enspace%
  \index{théorème de Weierstrass}\index{approximation uniforme}
  Si $f : [0,1] \to \mathbb{R}$ est continue, alors $\Bn(f)$ converge
  \emph{uniformément} vers $f$ sur $[0,1]$ :
  \[
    \bigl\| \Bn(f) - f \bigr\|_\infty
    \;=\; \sup_{x \in [0,1]} \bigl| \Bn(f)(x) - f(x) \bigr|
    \;\xrightarrow[\;n \to +\infty\;]{}\; 0.
  \]
\end{tcolorbox}

\begin{significationintuitive}
La preuve de Weierstrass par les polynômes de Bernstein est
\emph{constructive} : elle exhibe une suite explicite de polynômes
approchant $f$, là où l'énoncé classique se contente d'en affirmer
l'existence. Son ressort est \emph{probabiliste} : on lance $n$ fois une
pièce de probabilité $x$ et l'on évalue $f$ à la fréquence observée de
« piles ». La quantité $\Bn(f)(x)$ n'est autre que la \emph{moyenne pondérée}
des valeurs $f(i/n)$, les poids $\Bin(x)$ étant ceux d'une loi binomiale.
Par la loi des grands nombres, cette fréquence se concentre autour de $x$,
si bien que la moyenne pondérée de $f$ tend vers $f(x)$.
\end{significationintuitive}

\decsep{}

%=============================================================
%  § 2 — DÉMONSTRATION
%=============================================================
\secnum{navyblue}{2}{Démonstration}

\begin{ideepreuve}
Soit $S_n \sim \mathcal{B}(n,x)$ une variable binomiale : $\mathbb{P}(S_n=i)
= \Bin(x)$. Alors $\Bn(f)(x) = \mathbb{E}\!\left[f(S_n/n)\right]$. Comme
$\mathbb{E}[S_n/n]=x$ et $\mathrm{Var}(S_n/n)=x(1-x)/n \to 0$, la fréquence
$S_n/n$ se \emph{concentre} autour de $x$ (Bienaymé--Tchebychev). La
continuité uniforme de $f$ transforme cette concentration en convergence
de $\mathbb{E}[f(S_n/n)]$ vers $f(x)$, et l'uniformité en $x$ provient des
bornes uniformes fournies par le théorème de Heine.
\end{ideepreuve}

\vspace{10pt}
\rubriq{Preuve}

\begin{mdframed}[style=proofbar]
  \textit{Étape 1 — Les trois moments.}\enspace%
  \index{inégalité de Bienaymé-Tchebychev}
  La formule du binôme $\sum_{i=0}^n \binom{n}{i}x^i{(1-x)}^{n-i}={(x+(1-x))}^n=1$
  donne la partition de l'unité. En dérivant l'identité
  ${(x+y)}^n=\sum\binom{n}{i}x^i y^{n-i}$ par rapport à $x$ (puis $y=1-x$),
  on obtient successivement
  \[
    \sum_{i=0}^{n} \Bin(x) = 1, \qquad
    \sum_{i=0}^{n} \frac{i}{n}\,\Bin(x) = x, \qquad
    \sum_{i=0}^{n} {\left(\frac{i}{n}-x\right)}^{\!2} \Bin(x)
       = \frac{x(1-x)}{n}.
  \]
  Les deux premières expriment $\Bn(1)=1$ et $\Bn(\mathrm{id})=\mathrm{id}$.

  \medskip
  \textit{Étape 2 — Continuité uniforme (Heine).}
  $f$ étant continue sur le compact $[0,1]$, elle y est bornée, disons
  $|f|\leq M$, et \emph{uniformément continue} : pour $\varepsilon>0$ donné,
  il existe $\delta>0$ tel que
  \[
    |u-v|\leq\delta \;\Longrightarrow\; |f(u)-f(v)|\leq\varepsilon.
  \]

  \medskip
  \textit{Étape 3 — Découpage selon $|i/n-x|$.}
  Fixons $x\in[0,1]$. En utilisant $\sum_i\Bin(x)=1$ :
  \[
    \bigl|\Bn(f)(x)-f(x)\bigr|
    \;\leq\; \sum_{i=0}^{n}\left|f\!\left(\tfrac{i}{n}\right)-f(x)\right|\Bin(x).
  \]
  On scinde la somme selon que $\bigl|\tfrac{i}{n}-x\bigr|\leq\delta$ ($A$)
  ou $>\delta$ ($B$). Sur $A$, l'uniforme continuité donne
  $\sum_{A}\varepsilon\,\Bin(x)\leq\varepsilon$. Sur $B$, on majore par
  $2M$ et l'on utilise $1 \leq \tfrac{1}{\delta^2}{\left(\tfrac{i}{n}-x\right)}^2$ :
  \[
    \sum_{B}2M\,\Bin(x)
    \;\leq\; \frac{2M}{\delta^2}\sum_{i=0}^{n}
             {\left(\tfrac{i}{n}-x\right)}^2\Bin(x)
    \;=\; \frac{2M}{\delta^2}\cdot\frac{x(1-x)}{n}.
  \]

  \medskip
  \textit{Étape 4 — Majoration finale uniforme.}
  Comme $x(1-x)\leq\tfrac14$ sur $[0,1]$, on obtient, \emph{indépendamment
  de $x$} :
  \[
    \bigl\|\Bn(f)-f\bigr\|_\infty
    \;\leq\; \varepsilon + \frac{M}{2\,\delta^2\,n}.
  \]
  Pour $n\geq M/(2\delta^2\varepsilon)$, le majorant est $\leq 2\varepsilon$.
  Comme $\varepsilon>0$ est arbitraire, $\|\Bn(f)-f\|_\infty\to 0$.
  \hfill$\blacksquare$
\end{mdframed}

\decsep{}

%=============================================================
%  § 3 — EXEMPLES, CONTRE-EXEMPLES, EXERCICES
%=============================================================
\secnum{exgreen}{3}{Exemples, contre-exemples et exercices}

\rubriq{Exemples}

\begin{mdframed}[style=exbar]
  \textbf{1. Courbes de Bézier.}\enspace\index{courbes de Bézier}
  En conception assistée par ordinateur (CAO), la courbe de Bézier de points
  de contrôle $P_0,\dots,P_n\in\mathbb{R}^2$ est
  \[
    \gamma(t) = \sum_{i=0}^{n} P_i\,B_{i,n}(t),\qquad t\in[0,1].
  \]
  La partition de l'unité et la positivité des $B_{i,n}$ assurent que
  $\gamma(t)$ reste dans l'enveloppe convexe des $P_i$.

  \smallskip\noindent
  \textbf{2. Image d'une fonction simple.}
  Pour $f(x)=x^2$, un calcul direct (ou les trois moments) donne
  \[
    \Bn(f)(x) = \frac{n-1}{n}\,x^2 + \frac{1}{n}\,x
             = x^2 + \frac{x(1-x)}{n}.
  \]
  On lit l'erreur $\Bn(f)-f = x(1-x)/n \to 0$ uniformément en $1/n$.

  \smallskip\noindent
  \textbf{3. Lien avec la loi binomiale.}\enspace%
  $\Bn(f)(x)=\mathbb{E}\!\left[f(S_n/n)\right]$ avec $S_n\sim\mathcal{B}(n,x)$ ;
  les poids $\Bin(x)$ sont exactement les probabilités $\mathbb{P}(S_n=i)$.
\end{mdframed}

\vspace{6pt}
\rubriq{Contre-exemples et pièges}

\begin{mdframed}[style=cexbar]
  \textbf{Convergence lente.}\enspace
  Pour $f$ régulière, $\Bn(f)-f = O(1/n)$ seulement (cf. $f(x)=x^2$ ci-dessus) :
  Bernstein n'est \emph{pas} fait pour l'interpolation rapide. Les nœuds de
  Tchebychev offrent une convergence bien plus véloce pour les fonctions
  lisses.

  \smallskip\noindent
  \textbf{Ce n'est pas de l'interpolation.}\enspace
  En général $\Bn(f)$ ne passe \emph{pas} par les points $(i/n,\,f(i/n))$,
  sauf aux extrémités $x=0$ et $x=1$. C'est de l'\emph{approximation},
  non de l'interpolation : $\Bn(f)$ « lisse » $f$ au lieu de l'épouser.

  \smallskip\noindent
  \textbf{Ne pas confondre avec Lagrange.}\enspace
  Le polynôme d'interpolation de Lagrange passe par tous les nœuds mais peut
  \emph{diverger} (phénomène de Runge) ; $\Bn(f)$ converge toujours
  uniformément, au prix d'une vitesse modeste.
\end{mdframed}

\vspace{6pt}
\exolabel

\begin{mdframed}[style=exobar]
  \begin{enumerate}
    \item \textbf{Les trois moments.} Établir
          $\sum_i \Bin(x)=1$, $\sum_i \tfrac{i}{n}\Bin(x)=x$ et
          $\sum_i {\left(\tfrac{i}{n}-x\right)}^2\Bin(x)=\tfrac{x(1-x)}{n}$
          en dérivant ${(x+y)}^n$ par rapport à $x$.

    \item \textbf{Image d'un polynôme.} Montrer que si $f$ est un polynôme de
          degré $d$, alors $\Bn(f)$ est un polynôme de degré $\leq d$, et que
          $\Bn(f)\to f$ converge en réalité en $O(1/n)$.

    \item \textbf{Reproduction des affines.} Déduire des moments que
          $\Bn(ax+b)=ax+b$ pour tous $a,b\in\mathbb{R}$.

    \item \textbf{Dérivée.} Établir la formule
          \[
            \Bn'(f)(x) = n\sum_{i=0}^{n-1}
              \left[f\!\left(\tfrac{i+1}{n}\right)-f\!\left(\tfrac{i}{n}\right)\right]
              B_{i,n-1}(x),
          \]
          et en déduire que $\Bn(f)$ est croissante dès que $f$ l'est.

    \item \textbf{Module de continuité.} Avec $\omega_f(\delta)
          =\sup_{|u-v|\leq\delta}|f(u)-f(v)|$, montrer
          $\|\Bn(f)-f\|_\infty \leq \tfrac{3}{2}\,\omega_f\!\bigl(1/\sqrt{n}\bigr)$.

    \item \textbf{Cas lipschitzien.} Si $f$ est $k$-lipschitzienne, montrer,
          en majorant directement $\sum_i \bigl|\tfrac{i}{n}-x\bigr|\,\Bin(x)$
          par l'inégalité de Cauchy--Schwarz, que
          $\|\Bn(f)-f\|_\infty \leq \tfrac{k}{2\sqrt{n}}$.
  \end{enumerate}
\end{mdframed}

\leconfooter{S.\ Bernstein (1912)}
